\begin{textsample}{NEG Dim 6 – global\_north\_2024\_04 – Score 1.00 – global\_north\_2024\_04/200442352.txt}  \label{ex:f6_neg_374}
Netanyahu Says He ' ll Invade Rafah"With or Without a Deal"Israeli Prime Minister Benjamin Netanyahu speaks during a joint press conference with the German Chancellor after their meeting in Jerusalem on March 17, 2024.
LEO CORREA/Getty Images The United States spent Monday pushing Hamas to accept the"generous"new ceasefire proposal Israel offered up, with Secretary of State Antony Blinken suggesting the deal could head off Prime Minister Benjamin Netanyahu ' s plan to expand his military operation into Rafah.
But on Tuesday, Netanyahu seemed to undermine the U.S. line, vowing to forge ahead with a Rafah invasion whether Hamas accepts the truce plan or not."The idea that we will stop the war before achieving all of its goals is out of the question,"the prime minister said in a statement."We will enter Rafah and we will eliminate Hamas ' battalions there -- with or without a deal, to achieve the total victory."The pledge could dim the brief glimmer of hope that has emerged in recent days after more than seven months @ @ @ @ @ @ @ @ @ @.
For weeks, negotiations between Israel and Hamas have been deadlocked, despite international pressure for Hamas to release its hostages and Israel to deescalate its bombing of Gaza.
But over the weekend, Israel softened its demands, reportedly calling for the release of 33 hostages in the first stage of the deal, with the possibility of the"restoration of sustainable calm"in Gaza afterward."It is possible that if the first stage is implemented,"an Israeli official told Axios,"it will be possible to advance to the next stages and reach the end of the war."Hamas didn't immediately issue a formal comment on the proposal, but an official with the group suggested it had"no major issues"with it, and the U.S. called for the organization to accept the deal quickly :"The only thing standing between the people of Gaza and a ceasefire is Hamas,"Blinken said in Saudi Arabia on Monday."I ' m hopeful they will make the right decision. @ @ @ @ @ @ @ @ @ @ again threatening to defy President Joe Biden, who has warned that an incursion into Rafah -- where hundreds of thousands of civilians are sheltering -- would worsen the humanitarian crisis in Gaza and cross a"red line."Advocates of the truce deal had hoped Hamas would be incentivized to accept Israel ' s new terms if it could prevent the Rafah operation.
But by suggesting the deal will not have any bearing on his Rafah plans, Netanyahu may be dooming its prospects.
That could both exacerbate the crisis in Gaza, where famine appears imminent, and further complicate the effort to reclaim hostages from Hamas.
It also threatens to increase tensions with Biden, a steadfast supporter of Israel who is facing mounting pressure at home over the war as campus protests escalate -- a potential preview of the unrest he could see at the Democratic National Convention in Chicago this summer, just months before his November rematch against Donald Trump. ( The former president, in a Time interview published Tuesday, said Netanyahu bears responsibility for the @ @ @ @ @ @ @ @ @ @.") As Haaretz ' s Amos Harelpointed out Tuesday, Netanyahu ' s defiance seems even to put him at odds with much of his own war cabinet,"most of whose members now prefer to promote a deal at the expense of invading Rafah."But Bibi seems to be siding with his allies on the far-right, who are openly seeking to implode the U.S.-backed deal :"The prime minister heard my words,"Itamar Ben Gvir, Israel ' s far-right national security minister, said after meeting with Netanyahu, and"promised that the war would not end, and promised that there would be no reckless deal."

% matched lemmas: 
\end{textsample}
